\documentclass[b5j,fleqn]{jarticle}
\usepackage{zennyakuUniv}
\begin{document}
\body{
\Title{newcrown3 p47}
\para

\sentP{Dinu}{This movie is in Hindi, but it has English subtitles.}{ディヌ}{この映画はヒンディー語ですが、英語の字幕が付いています。}

\sentP{Riku}{Oh, good.}{陸}{ああ、よかった。}

\sentP{Riku}{Are all Indian movies in Hindi?}{陸}{インド映画はすべてヒンディー語なのですか。}

\sentP{Dinu}{No.}{ディヌ}{いいえ。}

\sentP{Dinu}{About half are in Hindi.}{ディヌ}{約半分がヒンディー語です。}

\sentP{Dinu}{However, many languages are spoken in India.}{ディヌ}{しかし、インドでは多くの言語が話されています。}

\sentP{Dinu}{So there are movies made in other languages like Tamil or Telugu.}{ディヌ}{そのため、タミル語やテルグ語のような、ほかの言語で作られた映画もあります。}

\sentP{Riku}{Wow, how do you watch those movies?}{陸}{へえ、そうした映画をどのように見るのですか。}

\sentP{Riku}{Do you set subtitles?}{陸}{字幕を設定するのですか。}

\sentP{Dinu}{Yes.}{ディヌ}{はい。}

\sentP{Dinu}{This streaming service is very convenient because I can choose subtitles in Hindi or English.}{ディヌ}{このストリーミングサービスでは、ヒンディー語か英語の字幕を選べるので、とても便利です。}
}
\end{document}
